\usetikzlibrary{decorations.markings, arrows.meta} % Untuk tanda segmen yang sama panjang

% Definisi untuk tanda segmen yang sama panjang (seperti di sketsa)
\tikzset{
    seg mark/.style={
        decoration={markings, mark=at position 0.5 with {\draw[#1] (-2pt,-2pt) -- (2pt,2pt) (2pt,-2pt) -- (-2pt,2pt);}},
        postaction={decorate}
    },
    seg mark half/.style={
        decoration={markings, mark=at position 0.5 with {\draw[#1] (0pt,-2pt) -- (0pt,2pt);}},
        postaction={decorate}
    },
    seg mark double/.style={
        decoration={markings, mark=at position 0.5 with {\draw[#1] (-2pt,-2pt) -- (-2pt,2pt) (2pt,-2pt) -- (2pt,2pt);}},
        postaction={decorate}
    },
    seg mark triple/.style={
        decoration={markings, mark=at position 0.5 with {\draw[#1] (-4pt,-2pt) -- (-2pt,2pt) (0pt,-2pt) -- (0pt,2pt) (2pt,-2pt) -- (4pt,2pt);}},
        postaction={decorate}
    },
}

\begin{center}
\begin{tikzpicture}[scale=1.5]
    % 1. Definisikan segitiga sama kaki ABC (AB=AC)
    \tkzDefPoint(0,0){B}
    \tkzDefPoint(4,0){C}
    \tkzDefPoint(2,4){A} % Titik A di tengah-tengah atas BC
    
    % Gambar segitiga ABC
    \tkzDrawPolygon(A,B,C)

    % 2. D adalah titik tengah AC, E adalah titik tengah BC
    \tkzDefMidPoint(A,C)\tkzGetPoint{D}
    \tkzDefMidPoint(B,C)\tkzGetPoint{E}
    
    % 3. M adalah titik tengah DE
    \tkzDefMidPoint(D,E)\tkzGetPoint{M}
    
    % 4. F adalah perpotongan AM dan BC
    \tkzInterLL(A,M)(B,C)\tkzGetPoint{F}
    
    % --- GAMBAR OBJEK UTAMA ---
    
    % Garis-garis yang menghubungkan titik-titik konstruksi
    \tkzDrawSegments(A,D D,E A,M D,F M,F)
    
    % Gambar lingkaran luar ADM (yang akan disentuh oleh DF)
    \tkzDefCircumCenter(A,D,M)\tkzGetPoint{O_ADM}
    \tkzDrawCircle[blue](O_ADM,A) % Gambar lingkaran luar ADM
    
    % --- TANDA SEGMEN SAMA PANJANG ---
    \tkzDrawSegments[seg mark=black](A,D D,C) % D midpoint AC
    \tkzDrawSegments[seg mark half=black](B,E E,C) % E midpoint BC
    \tkzDrawSegments[seg mark triple=black](D,M M,E) % M midpoint DE

    % Sudut
    \pic [fill=blue!20, angle eccentricity=1.5] {angle = F--D--C};
    \pic [fill=blue!20, angle eccentricity=1.5] {angle = B--A--F};
    \pic [draw=red, ->, angle radius=20pt] {angle = F--A--D};
    \pic [draw=red, ->, angle radius=20pt] {angle = M--D--F};

    % --- LABEL TITIK ---
    \tkzDrawPoints[fill=black, size=3](A,B,C,D,E,M,F)
    \tkzLabelPoints[above](A)
    \tkzLabelPoints[below left](B)
    \tkzLabelPoints[below right](C)
    \tkzLabelPoints[right](D)
    \tkzLabelPoints[below](E,F)
    \tkzLabelPoints[left](M)
\end{tikzpicture}
\end{center}